\chapter{Guia de Uso}

\begin{planningbox}{Como esta construido este suplemento}
Este cuaderno no repite la practica basica del libro principal. Cada problema parte de una
tipologia ya trabajada y la traslada a un contexto donde hace falta:
\begin{itemize}
  \item reconocer que tecnica manda;
  \item traducir datos a lenguaje matematico;
  \item sostener una representacion visual;
  \item interpretar el resultado final.
\end{itemize}
\end{planningbox}

\begin{databox}{Alcance de esta version}
La presente iteracion deja desarrollado el bloque \texttt{C05-C06}, con un problema largo por
seccion, respuesta breve y solucion completa. El apendice final mantiene visible el mapa de
todo el suplemento para continuar la expansion por capitulos.
\end{databox}

\begin{methodbox}{Modo recomendado de trabajo}
\begin{enumerate}
  \item Lee primero el contexto y localiza magnitudes, unidades y relaciones.
  \item Decide que parte del temario domina la situacion.
  \item Haz un dibujo, tabla o esquema antes de operar.
  \item Resuelve con forma exacta cuando sea razonable.
  \item Cierra con una frase de interpretacion, no solo con un numero.
\end{enumerate}
\end{methodbox}
